
\lussec 10. 结束语

通过向夸克场的推广，
梯度流成为研究 QCD 中手征对称性自发破缺
以及轻赝标介子物理的有力工具。
本文的重点是理论框架及其在广泛使用的格点 QCD 表述之一中的实现。
鉴于流的重正化性质，
可以预期：只要格点理论的局域性与规范不变性得到保证，
格点正规化的选择在连续极限下无关紧要。
显然，这一理论分析同样适用于
具有其他规范群与费米子多重态的场论。

在迄今具体完成的例子中，
梯度流在数值格点 QCD 中的应用
相对于既有的计算策略展现出重要的技术优势。
除了第 1 节已经提到的那些之外，
还存在许多潜在的进一步用途。
例如味单态道尚未被考虑，
而随时间演化的凝聚很可能是非零温度 QCD 中
一个易于获取的手征序参量。
此外，把流应用于共形窗附近的类 QCD 理论，
或许可以带来有趣的定性洞见。
